% SGLang, TokenFlow and TensorRT-LLM: owner, 2026-11-05; which SGLang algorithms and where are not on record, so no Fraunhofer bullet or tag.
% GStreamer, WebRTC and Jetson: industry projects (owner, 2026-11-03); safety requirements met by batching (owner, 2026-10-07), no figures on record.
\IfLanguageName{english}{
    \par{At Fraunhofer IPA I serve a multimodal model on vLLM for our coding agent, so long contexts and tool calls are what I tune for. For vLLM I reimplemented algorithms from SGLang and TokenFlow (EuroSys 2026), and I have used \mbox{TensorRT-LLM}. In industry projects I ran GStreamer camera pipelines with WebRTC and Rust/C++ inference and live detection on a Jetson GPU. For a customer's safety system I met strict latency and throughput requirements by batching 4 cameras.}
}{
    \par{Am Fraunhofer IPA betreibe ich ein multimodales Modell auf vLLM für unseren Coding-Agenten, lange Kontexte und Tool-Calls sind also die Last, für die ich optimiere. Für vLLM habe ich Algorithmen aus SGLang und TokenFlow (EuroSys 2026) neu implementiert und TensorRT-LLM eingesetzt. In Industrieprojekten habe ich GStreamer-Kamera-Pipelines mit WebRTC und Rust/C++-Inferenz sowie Live-Objekterkennung auf einer Jetson-GPU betrieben. Für das Sicherheitssystem eines Kunden habe ich strenge Latenz- und Durchsatzanforderungen mit Batch-Inferenz über 4 Kameras erfüllt.}
}
